% !TeX root = ../papers/nc-p3.tex
\section{二十一次母線束上での全隣接辺と整余因子曲線の排除}
\label{sec:bezout-scroll-neighbors}
\begin{theorem}\label{thm:exclude-integral-cofactor-curves}
標数零の代数閉体上の標準quadratic Koszul左右AS正則な四辺有限被約窓では、
非恒等な反対平方根の余因子スキームの正次元部分全体が、整な台を持つ一般被約曲線となることはない。
\end{theorem}
\begin{proof}[母線束への任意の辺の引戻し]
残る二十一次では入力 $x\in\mathbb P(S_t^3)$ と反対 $y\in\mathbb P(S_s^3)$ は
$f_x=T(t_0x),g_x=T(t_1x)$ の三次共通因子により対応する。
$s_0,s_1$ と $t_0,t_1$ は二つの二項一次基底である。
正規化 $\nu:S\to D$ は $\mathbb P^1_s\times\mathbb P^3_x$ 内で
$f_x(z)=g_x(z)=0$ により定まり、射影直線束 $\pi:S\to\mathbb P^1$ である。
超平面束 $H$ と母線 $F$ は $H^2=8,HF=1$。
三次共通因子の軌跡は階数二の整曲線上の真閉集合なので有限である。
標準の一次シジジー辺 $(x,C)$ はこの上に一意の $C$ を与え、両射影が有限である。
実際 $f_x,g_x$ は常に独立で、線形係数シジジーがあれば三次共通因子があり一意となる。
六双一次完全交叉の次数は二十で既存の被約二十印で尽きる。
四次共通因子はなく、各印での正規化繊維は $\mathbb P^1$ 内の長さ三スキームである。

任意の実際の隣接辺の六行四列一次行列を $E(x)$ とする。正規化上の余核の反射包を
$\mathcal V=(\operatorname{coker}(\mathcal O_S(-H)^4\to\mathcal O_S^6))^{**}$ と書く。
これは階数二のベクトル束であり、六切断が有限集合の外で生成し、
\[
                       c_1(\mathcal V)=4H,\qquad c_2(\mathcal V)=20 .
\]
実際階数落ちのイデアルの引戻しは二十印の繊維で、長さの和は六十。
各局所繊維は滑らかな曲線 $\mathbb P^1$ 内の因子なので、滑らかな曲面 $S$ 内で局所完全交叉である
（埋込み次元一以下の零次元局所環のイデアルはこの場合二生成）。
三つの単元を消去すれば残りは一列三成分で、この二生成イデアルを生成する。
局所基底変更により余核はイデアルと自由一元の和となり、反射包への余長がちょうど繊維長である。
最初の列写像は一般単射、従って単射であり、表示の第二Chern数 $10H^2=80$ から六十を引けばよい。

一般母線での分裂を $(a,4-a)$、$2\le a\le4$ とする。六切断による生成から両次数は非負である。
関数体上の一般母線で高い次数の直線部分束を選び、$S$ 上で飽和させる。
同次数でも一般母線は底関数体上の射影直線で、一本を選べる。飽和階数一部分層は反射的で、
商は直線束に有限スキームのイデアルを掛けたものである。従って
\[
 0\longrightarrow\mathcal O(aH+bF)\longrightarrow\mathcal V
   \longrightarrow I_Z(B)\longrightarrow0,\qquad B=(4-a)H-bF .
\]
$B$ の六切断には固定曲線がない。また
$20=8a(4-a)+(4-2a)b+\ell(Z)$。従って $a=2$ はなく、$a=3$ なら $b\ge2$、
$a=4$ なら $k=-b$ は $0\le k\le5$ である。特に $BH\le6$ である。
\end{proof}

\begin{proof}[商線束に現れる六独立切断]
上の六切断は独立である。従属なら、それらが定める有理写像 $S\dashrightarrow\mathbb P^5$ の像を含む
固定超平面 $R$ を取る。対応する点は $E(x)^t$ の核に入り $R$ にも入る。
六式のうち $R$ を基底に持つ五結合を行とする五行四列 $E_R$ を考える
（$R$ はここでは射影超平面、そのベクトル五空間を使う）。
五式のグラフの次元は一以下である。二次元以上の成分なら最後の豊富な $(1,1)$ 式を
加えても正次元基点が残るからである。従って $E_R$ の最大小行列式は余次元二の曲線を切り、
Hilbert--Burchで次数十である。一般平面では長さ十のスキーム $J$ を得る。
一般 $x\in D$ でこれらが全て零なら次数四の式が八次 $D$ を含んでしまう。
従って一般に $R$ 内の核点は唯一で、その座標はこれらの五つの四次式である。

一般平面切断の正規化である $H$ の一般入力切断へ制限する。
有限射のこの曲線上の繊維長は全て二以下である（平面は三次共通因子の有限集合を避ける）。
平面への有限押出しは各局所環上で高々二元生成であり、
$J$ のイデアルで割った商長の合計は高々二十となる。
従って五四次式の正規化上での共通固定次数は二十以下で、核点写像の超平面次数は
$4H^2-20=12$ 以上である。一方 $B$ の切断には固定曲線がなく、一般の入力切断はその孤立基点を避けるから
この次数は $BH\le6$ である。矛盾によって六切断は独立である。

$a=4$ なら $h^0(kF)\ge6$ より $k=5$。$a=3$ なら
$\pi_*\mathcal O(H)=\mathcal O(e)\oplus\mathcal O(8-e)$、$e\le4$ とし、
$B=H-bF$ は固定曲線がないので $b\le e$。よって $h^0(B)=10-2b\ge6$、
ゆえに $b=2$。両場合で $\ell(Z)=0$ となり、六切断はそれぞれ
$B=5F$ と $B=H-2F$ の全切断である。
\end{proof}

\begin{proof}[二つの線束からの列の復元]
$E$ の各列は、六切断への商を取ると
$H^0(\mathcal O_{\mathbb P^3}(1))\otimes H^0(B)\to H^0(H+B)$ の核に入る。
$B=5F$ なら入射の二式は $(4,1)$ なので、その $(5,1)$ 消滅式は
$s_i f_x,s_i g_x$ の四本で尽きる。実際Koszul分解の左端は $(-3,-1)$ で全コホモロジーが消える。
従って四列はこれらの基底となり $E$ は標準の $(x,C)$ 辺である。
実際の $(C,y)$ 辺も標準辺である。各印で標準二辺の関係が成立し、
逆変換 $T^{-1}$ でも共通三割平面は同じ有限集合（平面の両側の三次因子は一意）である。
従って標準辺は有限完全交叉で、そのグラフの双一次消滅式は六個の定義式で尽きる。
定理\ref{thm:trisecant-middle-simple-obstruction}に反する。

$B=H-2F$ の場合の列を明記する。$V_x$ を入力一次式の空間とする。
入射のKoszul分解を $(-2,1)$ 捻ると、底のSerre双対により
\[
 H^0(B)=\{(a,b)\in((S_s^4)^*)^2:
                  a(f_x)+b(g_x)=0\text{ for all }x\}.
\]
これは $\gamma\in(S_t^5)^*$ と
$a(l)=\gamma(t_1T^{-1}l),\ b(l)=-\gamma(t_0T^{-1}l)$ により同型である。
全射性は $S_t^4\oplus S_t^4\to S_t^5$, $(h,k)\mapsto t_1h-t_0k$
の核が $(t_0x,t_1x)$ であるためである。

次に $(-2,2)$ 捻りでは、$\mathbb P^1\times\mathbb P^3$ の三項の底次数は
$-10,-6,-2$、入力次数は $0,1,2$ である。全てで底の $H^1$ だけが残るので、
$H^0(2H-2F)$ は
\[
 (S_s^8)^* \longrightarrow ((S_s^4)^*)^2\otimes V_x
     \longrightarrow \operatorname{Sym}^2 V_x
\]
の中央コホモロジーである。第一写像は単射で、
$\eta\mapsto (\eta(g_x\,\cdot),-\eta(f_x\,\cdot))$ である
（単射性はKoszul分解の負全次数コホモロジー消滅からも従う）。
$V_x\otimes H^0(B)$ はこの中央に単射で、その積はコホモロジーへの自然写像である。
従って積の核は第一写像の像のうち各 $x$ で $H^0(B)$ の条件を満たすもの、
すなわち全 $x,x'$ について
\[
                \eta(g_x f_{x'}-f_x g_{x'})=0
\]
を満たす $\eta$ である。二項四次上のHankel形式 $G(l,h)=\eta(lh)$ で書くと
これは $T^t G T$ もHankelであることに等しい。実際四次単項式の行列で隣接する反対角成分が
同じという等式である。$G(Tk,Tk')=\psi_G(kk')$、$G(l,h)=\phi_G(lh)$ と書けば、
対応する列の $x$ での値は五次汎関数 $h\mapsto\psi_G(xh)$ である。
上記の $a,b$ に代入するとそれぞれ $\eta(g_x\,\cdot),-\eta(f_x\,\cdot)$ となる。

従ってこの中間位置 $U$ は同時Hankel形式のある四空間であり、$(x,U)$ 辺は $\psi_G(x S_t^5)=0$。
各印でこれは $t_0x,t_1x$ が $T^t G T$ の核に入ること、従って
$T(t_0x),T(t_1x)$ の張る $y S_s^1$ が $G$ の核に入ることを意味する。
これより $(U,y)$ の印は $\phi_G(y S_s^5)=0$ を満たす。この最後の六式も有限な完全交叉を定める。
実際二項八次汎関数の中央Hankel行列が階数三以下であることは、その三次収縮に非零核があることと同値である。
一方は三次核の二つの一次倍を取ればよい。逆は、三次核がなければ随伴消去環
$k[s,t]/\operatorname{Ann}(\eta)$ の三次と五次がともに四次元となることを使う。
四次に非零の消去式があれば、一般一次式でさらに割った商は四次以降で零なので
そのHilbert関数は四次以降増加できない。従って中央の値は少なくとも四である。
しかも中央階数三で三次核の次元が二以上ならその一次倍は少なくとも三次元となり矛盾する。
階数二以下なら三次収縮階数も二以下である（同じHilbert関数の単調性と対称性を使う）。
従って $U$ で可能な $G$ は既存の有限 $(x,U)$ 辺の二十印であり、どれも中央階数三、
左右どちらの二項構造でも三次核がただ一つである。$U,y$ の六式の零集合は有限で、
次数二十により元の被約グラフに等しく、その双一次式はこれらの六式で尽きる。

こうして $U$ の両隣は同じ四空間上の二つのHankel表示となり、その二つの四次評価曲線は異なる
（整な八次で二曲線は交わらない）。定理\ref{thm:exclude-degree-twenty-hankel-pairs}の
異なる構造の証明は、この二つの隣接辺と三次交叉次元四だけを使う。
よって残る中間位置が標準三割平面の基底変更行列となり、やはり
定理\ref{thm:trisecant-middle-simple-obstruction}が正則性を排除する。
以上で任意の隣接辺を尽くした。
\end{proof}
\path{computations/verify_bezout_scroll_neighbors.py} は二つの分裂型・残余長と、
Koszul接続写像からの同時Hankel列を検算する。
